% Aquí debes exponer las leyes de la física y los principios
% que rigen el comportamiento de la mezcla, independientemente de la sustancia
% computación que sustentan tu herramienta. 
% No defines términos, sino que explicas teorías.
\section{Marco Teórico}
\label{sec:marco_teorico}

    \subsection{Fundamentos de las Ecuaciones de Estado}
        Para \textcite{bib:thermo:goodstein2002states}, una ecuación de estado es toda relación entre la presión (\pressure[]{}) y temperatura (\temperature[]{}) de un sistema termodinámico en estado gaseoso.

        
    Justificación del uso de la correlación de Wagner para sustancias que no sean agua, representar el comportamiento cerca del punto crítico.



    \subsection{Termodinámica de Mezclas de Gases Ideales}
        El comportamiento de la mezcla de gases se basa el la ley de Dalton

    \subsection{Naturaleza del Equilibrio}
        Según \textcite{bib:thermo:VanNess}, el equilibrio es una condición estática, donde no ocurre cambio alguna en las propiedades macroscópicas de un sistema en el tiempo.
        Para \textcite{bib:thermo:goodstein2002states} un sistema en equilibrio, es cuando el estado inicial se vuelve irrelevante; lo que significa que, independientemente del mecanismo de cambio de estado, de la manera en que las partículas puedan redistribuir energía y cantidad de movimiento entre sí, todo recuerdo de cómo comenzó el sistema debe eventualmente quedar borrado por las infinitas distintas formas que las partículas puedan redistribuirse. Cuando un sistema alcanza esta condición, se dice que está en \textit{equilibrio}.

        y \textcite{bib:thermo:gibbs1874} dice: "Para el equilibrio de cualquier sistema aislado es necesario y suficiente que en todas las posibles variaciones de el estado del sistema que no alteran su energía interna la variación de su entropía sea menor o igual que cero"; es decir, un sistema está en equilibrio cuando ya no puede reorganizarse de ninguna manera que lo haga más desordenado.

    \subsection{Tipos del Equilibrio}
        Estos son algunas formas de equilibrio:
        \subsubsection{Equilibrio Térmico}
            Cuando la temperatura es uniforme en todo el sistema.
        \subsubsection{Equilibrio Mecánico}
            Cuando la presión es uniforme y no hay fuerzas que desequilibren el sistema
        \subsubsection{Equilibrio Químico}
            La composición química no cambia;
        \subsubsection{Equilibrio termodinámico}
            Cuando el sistema alcanza las tres formas de equilibrio (térmico, mecánico, químico).

        \subsection{Equilibrio de Fases}
            Una fase es una región de un sistema con propiedades químicas y físicas uniformes; dentro de un sistema pueden coexistir diversas fases. Las transiciones de una fase a otra están reguladas por los principios termodinámicos.

            \eqnClausiusClapeyron
    \subsection{Termodinámica del Equilibrio de Fases}

        % \genericPhaseDiagram

        \begin{figure}[h]
            \centering
            % --- DEFINICIONES PARA DIAGRAMA GLOBAL ---
\def\TminGlobal{-50}
\def\TmaxGlobal{450} % Supera los 373.9°C del punto crítico
\def\PminGlobal{0.1}
\def\PmaxGlobal{30500} % Supera los 22064 kPa del punto crítico

\begin{tikzpicture}
\begin{axis}[
    width=12cm, height=9cm,
    % Separación de los nombres de los ejes
    xlabel={Temperatura $\temperature[]{}$},
    xlabel style={yshift=-0.5cm},
    ylabel={Presión $\pressure[]{}$},
    ylabel style={xshift=-0.1cm},
    ylabel style={yshift=0.5cm},
    xmin=0, xmax=110,
    ymin=0, ymax=110,
    axis on top,
    axis lines=left,
    axis line style={line width=1.0pt},
    enlargelimits={abs=7pt, upper},
    ticks=none, 
    clip=false 
]
    % --- REGIONES (Fills corregidos) ---
    \fill[cyan!10] (axis cs:0,0) -- (axis cs:0,110) -- (axis cs:25,110) -- (axis cs:30,30) -- (axis cs:10,10) -- cycle;
    
    \fill[blue!10] (axis cs:30,30) -- (axis cs:25,110) -- (axis cs:80,110) -- (axis cs:80,75) 
        .. controls (axis cs:70,60) and (axis cs:50,45) .. (axis cs:30,30) -- cycle;
    
    \fill[orange!5] (axis cs:0,0) -- (axis cs:10,10) -- (axis cs:30,30) 
        .. controls (axis cs:50,45) and (axis cs:70,60) .. (axis cs:80,75) 
        -- (axis cs:110,75) -- (axis cs:110,0) -- cycle;
        
    \fill[purple!10] (axis cs:80,75) -- (axis cs:110,75) -- (axis cs:110,110) -- (axis cs:80,110) -- cycle;

    % --- CURVAS DE COEXISTENCIA ---
    \draw[teal, ultra thick] (axis cs:10,10) -- (axis cs:30,30);
    \draw[blue, ultra thick] (axis cs:30,30) -- (axis cs:25,110); 
    \draw[blue, ultra thick] (axis cs:30,30) .. controls (axis cs:50,45) and (axis cs:70,60) .. (axis cs:80,75);

    % --- PROYECCIONES Y LABELS EN EJES ---
    % Punto Triple
    \draw[gray, dotted, thick] (axis cs:30,30) -- (axis cs:30,0);
    \node[below, font=\small, yshift=-2pt] at (axis cs:30,0) {$\tTriple$};
    
    \draw[gray, dotted, thick] (axis cs:30,30) -- (axis cs:0,30);
    \node[left, font=\small, xshift=-2pt] at (axis cs:0,30) {$\pTriple$};
    
    % Punto Crítico
    \draw[gray, dashed, thick] (axis cs:80,75) -- (axis cs:80,0);
    \node[below, font=\small, yshift=-2pt] at (axis cs:80,0) {$\Tc$};
    
    \draw[gray, dashed, thick] (axis cs:80,75) -- (axis cs:0,75);
    \node[left, font=\small, xshift=-2pt] at (axis cs:0,75) {$\pC$};

    % --- PUNTOS FÍSICOS ---
    \filldraw (axis cs:30,30) circle (2.5pt) node[above left, font=\footnotesize, xshift=-2pt] {\textbf{P. Triple}};
    \filldraw (axis cs:80,75) circle (2.5pt) node[above right, font=\footnotesize] {\textbf{P. Crítico}};

    % --- ETIQUETAS DE ESTADO ---
    \node[blue!80!black] at (axis cs:12, 60) {\large \textbf{SÓLIDO}};
    \node[blue!70!black] at (axis cs:50, 85) {\large \textbf{LÍQUIDO}};
    \node[orange!70!black] at (axis cs:60, 20) {\large \textbf{VAPOR}};
    \node[purple!80!black, text width=2.5cm, align=center] at (axis cs:95, 92) {\footnotesize \textbf{Región Super Crítica}};

\end{axis}
\end{tikzpicture}
 % Llamada al segundo módulo
            
            \caption{Localización del punto de rocío en la línea de vapor saturado ($x=1$).}
        \end{figure}
        


\subsection{Fundamentos de la Termodinámica de Mezclas}
\begin{itemize}
    \item \textbf{Mezclas Gas-Vapor:} Comportamiento del aire seco y vapor de agua.
    \item \textbf{Ley de Dalton:} Presiones parciales en mezclas binarias.
    \item \textbf{Gases Ideales vs. Reales:} Justificación del modelo utilizado para el rango de operación.
\end{itemize}

\subsection{Teoría de la Psicrometría}
    \begin{itemize}
        \item \textbf{Propiedades Termodinámicas:} Definición matemática de Humedad Específica, Relativa y Entalpía.
        \item \textbf{Ecuaciones de Estado para el Vapor de Agua:} Modelos de presión de saturación (Ecuación de Antoine o Goff-Gratch).
        \item \textbf{Procesos Psicrométricos:} Calentamiento sensible, enfriamiento con deshumidificación y mezcla adiabática.
    \end{itemize}

    \subsection{Propiedades psicrométricas}
        \begin{itemize}
            \item Temperatura de bulbo seco
            \item Temperatura de bulbo húmedo
            \item Temperatura del punto de rocío
            \item Humedad relativa
            \item Relación de humedad
            \item Entalpía específica
            \item Volumen específico
        \end{itemize}




    \subsection{Entalpía específica del vapor}
        

\subsection{Geometría de la Carta Psicrométrica}
\begin{itemize}
    \item \textbf{Sistemas de Coordenadas:} Transformación de valores termodinámicos a coordenadas cartesianas ($x, y$).
    \item \textbf{Líneas de Saturación:} Deducción de la curva límite de la mezcla.
\end{itemize}



